%% kapa: exam=1 round=34 year=2023 period=1 form=A subject=economics subject_ko=경제학원론 seq=15 number=55
%% source: Q-Net 감정평가사 기출문제 (cst00309) artlSeq=5211789 file=raw/1008/5211789/2023년도 제34회 감정평가사 1차시험 1교시.pdf
%% conversion: kordoc 4.24.51 -> tools/gen1.py (auto); review=proofread-sonnet; flags=none
\begin{question}{55}
후생경제이론에 관한 설명으로 옳지 않은 것은?
\begin{choices}
  \item 계약곡선 위의 모든 점은 파레토효율적 배분을 대표한다.
  \item 일정한 전제하에서 왈라스균형은 일반경쟁균형이 될 수 있다.
  \item 차선의 이론에 따르면 점진적 접근방식에 의한 부분적 해결책이 최선은 아닐 수 있다.
  \item 후생경제학의 제1정리에 따르면 일반경쟁균형의 배분은 파레토효율적이다.
  \item 후생경제학의 제2정리는 재분배를 위한 목적으로 가격체계에 개입하는 것에 정당성을 부여한다.
\end{choices}
\end{question}
\answer{5}
